% !TEX root = task-1-15-1.tex
% Editable task bank. Append new BankTask blocks at the end of the document.
\documentclass[12pt,a4paper]{article}
\usepackage{physics-textbook}
\usepackage{graphicx,needspace}
\newif\ifShowSolutions
\ifdefined\TutorEdition
  \ShowSolutionstrue
\else
  \ShowSolutionsfalse % Change to \ShowSolutionstrue to display solutions.
\fi
\newcounter{banktask}
\newcommand{\BankTask}[3]{%
  \refstepcounter{banktask}%
  \Needspace{35mm}\par\medskip\noindent\begin{minipage}{\linewidth}%
  \noindent\textbf{\thebanktask.}\quad #2\par
  \ifShowSolutions
    #3\par
    \ifdefined\TutorEdition
      \ifcsname TutorSource#1\endcsname
        \smallskip{\small\csname TutorSource#1\endcsname\par}%
      \fi
    \fi
  \fi
  \end{minipage}\par\medskip
}
\begin{document}
\BookSubsection{1.15.1}{Работа}


% Source sdamgia-501: https://phys-ege.sdamgia.ru/problem?id=501
\BankTask{sdamgia-501}{
Ящик тянут по земле за веревку по горизонтальной окружности длиной $L=40\,\text{м}$ с постоянной по модулю скоростью. Работа силы тяги за один оборот по окружности $A=2{,}4\,\text{кДж}$. Чему равен модуль силы трения, действующей на ящик со стороны земли?
}{
\smallskip\noindent\textbf{Решение}\par
Поскольку ящик тянут с постоянной по модулю скоростью, его кинетическая энергия не меняется. Вся энергия, которая поступает в систему за счет работы силы тяги, расходуется на работу силы трения. Отсюда находим модуль силы трения:


$F_{\text{тр}}=\frac{A}{L}=\frac{2400\,\text{Дж}}{40\,\text{м}}=60\,\text{Н}.$

Ответ: $60$.
}

% Source sdamgia-529: https://phys-ege.sdamgia.ru/problem?id=529
\BankTask{sdamgia-529}{
Человек взялся за конец лежащего на земле однородного стержня длиной $2\,\text{м}$ и массой $100\,\text{кг}$ и поднял этот конец на высоту $1\,\text{м}$. Какую работу он совершил? Ускорение свободного падения принять равным $10\,\text{м/с}^2$.
}{
\smallskip\noindent\textbf{Решение}\par
Человек совершал работу против силы тяжести, то есть совершенная им работа шла на увеличение потенциальной энергии стержня. Поскольку он поднимал стержень за один конец, а второй конец оставался на земле, центр масс стержня в итоге оказался на высоте $0{,}5\,\text{м}$. Следовательно, потенциальная энергия стержня увеличилась на: 


$\Delta E_{\text{пот}}=mg\frac h2=100\,\text{кг}\cdot10\,\text{м/с}^2\cdot0{,}5\,\text{м}=500\,\text{Дж}.$

Именно этой величине и равна его работа.
 
Ответ: $500$.
}

% Source sdamgia-6335: https://phys-ege.sdamgia.ru/problem?id=6335
\BankTask{sdamgia-6335}{
Небольшое тело массой $200\,\text{г}$ свободно соскальзывает вниз по гладкой наклонной плоскости вдоль оси $Ox$. В таблице приведена зависимость проекции $\upsilon_x$ скорости этого тела от времени $t$. Какую работу совершит сила тяжести к моменту, к которому тело пройдет путь $1\,\text{м}$?
 
\begin{center}\begin{tabular}{|c|c|c|c|c|c|}\hline
$t$, с&$0$&$1$&$2$&$3$&$4$\\\hline
$\upsilon_x$, м/с&$0$&$0{,}5$&$1$&$1{,}5$&$2$\\\hline
\end{tabular}\end{center}
}{
\smallskip\noindent\textbf{Решение}\par
Из таблицы видно, что движение тела равноускоренное, следовательно, за время $t$ тело вдоль оси $OX$ проходит путь, равный $S=\frac{a_xt^2}{2}$ Найдем ускорение из таблицы: для этого нужно выбрать произвольный промежуток времени, вычесть из конечной скорости на промежутке начальную и разделить получившуюся разность на длительность промежутка времени $a=\frac{v_2-v_1}{t_2-t_1}$ Ускорение $a=\frac{0{,}5\,\text{м/с}-0\,\text{м/с}}{1\,\text{с}-0\,\text{с}}=0{,}5\,\text{м/с}^2$ Найдем, через какое время тело пройдет путь $1\,\text{м}$:


$S=\frac{at^2}{2}\ \Longleftrightarrow\ t=\sqrt{\frac{2S}{a_x}}$
 
$t=\sqrt{\frac{2\cdot1}{0{,}5}}=2\,\text{с}.$

Трения нет, поэтому выполняется закон сохранения механической энергии. Тело движется вдоль оси $OX$, при его скольжении вниз потенциальная энергия превращается в кинетическую. Работа, совершенная полем тяжести над телом, равна изменению потенциальной энергии взятой с противоположным знаком. получаем:


$E_{\text{к}1}+E_{\text{п}1}=E_{\text{к}2}+E_{\text{п}2}$
 
$A=-(E_{\text{п}2}-E_{\text{п}1})=E_{\text{к}2}-E_{\text{к}1}=\frac{mv_{x2}^2}{2}-\frac{mv_{x1}^2}{2}=\frac{0{,}2\,\text{кг}\cdot(1\,\text{м/с})^2}{2}=0{,}1\,\text{Дж}.$


Ответ: $0{,}1$.
}

% Source sdamgia-7344: https://phys-ege.sdamgia.ru/problem?id=7344
\BankTask{sdamgia-7344}{
Телу массой $2\,\text{кг}$, находящемуся у основания шероховатой наклонной плоскости, сообщили начальную скорость $3\,\text{м/с}$ в направлении вверх вдоль наклонной плоскости. Через некоторое время тело вернулось в исходную точку, имея втрое меньшую кинетическую энергию. Какую работу совершила сила трения за время движения тела?
}{
\smallskip\noindent\textbf{Решение}\par
По теореме о кинетической энергии изменение кинетической энергии равно работе внешних сил. Поскольку работа силы тяжести равна нулю, единственной силой, совершающей ненулевую работу, является сила трения. Тогда: 


$A_{\text{тр}}=\Delta E=E_{\text{кин}2}-E_{\text{кин}1}=\frac{E_{\text{кин}1}}3-E_{\text{кин}1}=-\frac23E_{\text{кин}1}=-\frac{2mV_1^2}{3\cdot2}=-6\,\text{Дж}.$


Ответ: $-6$.
}

% Source sdamgia-9140: https://phys-ege.sdamgia.ru/problem?id=9140
\BankTask{sdamgia-9140}{
К телу массой $5\,\text{кг}$, покоящемуся на шероховатой горизонтальной плоскости, в момент времени $t=0$ прикладывают горизонтально направленную силу $5\,\text{Н}$. Коэффициент трения между поверхностью тела и плоскостью равен $0{,}2$. Чему равна работа, совершаемая этой силой за первые $10$ минут ее действия?
}{
\smallskip\noindent\textbf{Решение}\par
Работа, совершаемая силой, равна $A=Fs\cos\alpha$ Чтобы тело можно было сдвинуть с места, должно выполняться условие: $F\ge F_{\text{тр}}$ Максимальная сила трения покоя равна:


$\mu mg=0{,}2\cdot5\cdot10=10\,\text{Н}.$

Прикладываемая сила меньше $10\,\text{Н}$, значит, тело не будет сдвинуто с места. Следовательно, работа будет равна $0$.
 
Ответ: $0$.
}

% Source sdamgia-29785: https://phys-ege.sdamgia.ru/problem?id=29785
\BankTask{sdamgia-29785}{
Тяжелый ящик тянут по горизонтальному полу со скоростью $0{,}3\,\text{м/с}$ при помощи троса, модуль силы натяжения которого постоянен и равен $160\,\text{Н}$. Трос все время составляет с полом угол $60^\circ$. Какую работу совершит сила натяжения троса за $10$ секунд, если ящик не отрывается от пола?
}{
\smallskip\noindent\textbf{Решение}\par
Работа силы натяжения равна $A=Fs\cos\alpha$ При прямолинейном равномерном движении $s=vt$ откуда $A=Fvt\cos\alpha=160\cdot0{,}3\cdot10\cdot0{,}5=240\,\text{Дж}.$
 
Ответ: $240$.
}

% Source sdamgia-37557: https://phys-ege.sdamgia.ru/problem?id=37557
\BankTask{sdamgia-37557}{
Тело массой $0{,}1\,\text{кг}$, подвешенное на нити длиной $1\,\text{м}$, вращается, двигаясь в горизонтальной плоскости по окружности. Чему равна работа, которую совершает сила тяжести, действующая на это тело, за время его полного оборота?
}{
\smallskip\noindent\textbf{Решение}\par
Работа силы тяжести по замкнутой траектории равна $0$.
 
Ответ: $0$.
}

% Source sdamgia-3260: https://phys-ege.sdamgia.ru/problem?id=3260
\BankTask{sdamgia-3260}{
Пуля массой $10\,\text{г}$, летящая со скоростью $200\,\text{м/с}$, пробивает доску толщиной $2\,\text{см}$ и вылетает со скоростью $100\,\text{м/с}$. Определите силу сопротивления доски, считая ее постоянной.
}{
\smallskip\noindent\textbf{Решение}\par
Запишем теорему о кинетической энергии: $E_{\text{к}2}-E_{\text{к}1}=A_{\text{сопр}}$ где начальная кинетическая энергия $E_{\text{к}1}=\frac{mv_1^2}{2}$ конечная кинетическая энергия $E_{\text{к}2}=\frac{mv_2^2}{2}$ работа силы сопротивления $A=F_{\text{сопр}}s\cos\alpha$ Перемещение пули в доске равно ее толщине $s=d$ угол между направлением вектора силы сопротивления и вектором перемещения равен $180^\circ$, поэтому $\cos\alpha=-1$ Следовательно, получаем: $-F_{\text{сопр}}d=\frac{mv_2^2}{2}-\frac{mv_1^2}{2}$ Таким образом, сила сопротивления доски равна


$F_{\text{сопр}}=\frac{m(v_1^2-v_2^2)}{2d}=\frac{0{,}01\,\text{кг}\bigl((200\,\text{м/с})^2-(100\,\text{м/с})^2\bigr)}{2\cdot0{,}02\,\text{м}}=7500\,\text{Н}.$

Ответ: $7500\,\text{Н}$.
}

% Source sdamgia-4430: https://phys-ege.sdamgia.ru/problem?id=4430
\BankTask{sdamgia-4430}{
Груз начинает свободно падать с некоторой высоты без начальной скорости. Пролетев $40\,\text{м}$, груз приобрел скорость $20\,\text{м/с}$. Чему, на этом участке пути, равно отношение изменения потенциальной энергии груза к работе силы сопротивления воздуха?
}{
\smallskip\noindent\textbf{Решение}\par
Запишем закон сохранения энергии для незамкнутой системы:


$A_{\text{сопр}}=E_{\text{пол}2}-E_{\text{пол}1}$

где полная энергия в верхней точке $E_{\text{пол}1}=E_{\text{п}}=mgh$ полная энергия в нижней точке $E_{\text{пол}2}=E_{\text{к}}=\frac{mv^2}{2}$
Тогда получаем:


$A_{\text{сопр}}=\frac{mv^2}{2}-mgh$ 

Изменение потенциальной энергии груза равно $\Delta E_{\text{п}}=0-mgh=-mgh$ 
Таким образом, отношение изменения потенциальной энергии груза к работе силы сопротивления воздуха равно: 


$\frac{\Delta E_{\text{п}}}{A_{\text{сопр}}}=\frac{-mgh}{-mgh+\frac{mv^2}{2}}=\frac{2gh}{2gh-v^2}=\frac{2\cdot10\,\text{м/с}^2\cdot40\,\text{м}}{2\cdot10\,\text{м/с}^2\cdot40\,\text{м}-(20\,\text{м/с})^2}=2.$


Ответ: $2$.
}

% Source sdamgia-10655: https://phys-ege.sdamgia.ru/problem?id=10655
\BankTask{sdamgia-10655}{
Брусок равномерно двигают по горизонтальной поверхности. В процессе движения по этой поверхности он проходит два участка одинаковой длины с различными коэффициентами трения. Известно, что на первом участке модуль работы силы трения $2{,}5\,\text{Дж}$, а на втором участке модуль работы силы трения $7{,}5\,\text{Дж}$. При этом на втором участке коэффициент трения на $0{,}4$ больше, чем на первом. Определите, чему равен коэффициент трения на втором участке.
}{
\smallskip\noindent\textbf{Решение}\par
При движении бруска по шероховатой поверхности будет возникать сила трения, равная:


$F_{\text{тр}}=\mu N$ где $N$ – сила реакции опоры.

Модуль работы силы трения равен:


$|A|=F_{\text{тр}}l$

Найдем отношение работ сил трения при прохождении двух различных участков:


$\frac{A_2}{A_1}=\frac{\mu_2Nl}{\mu_1Nl}=\frac{\mu_2}{\mu_1}=\frac{7{,}5}{2{,}5}=3$

По условию $\mu_1=\mu_2-0{,}4$. Тогда $3\mu_2-\mu_2=3\cdot0{,}4$, откуда $\mu_2=0{,}6$.

Ответ: $0{,}6$.
}


% Source 225EC6: https://ege.fipi.ru/bank/questions.php?search=1&proj=BA1F39653304A5B041B656915DC36B38&qid=225EC6
\BankTask{225EC6}{
При равномерном перемещении саней по горизонтальному участку пути длиной $50\,\text{м}$ сила тяги совершает работу, величина которой равна $1000\,\text{Дж}$. Какова сила трения? Выразите ответ в ньютонах.
}{
\smallskip\noindent\textbf{Решение от GPT}\par
При равномерном движении $\Delta E_{\text{к}}=0$, поэтому $A_{\text{тяги}}+A_{\text{тр}}=0$. На горизонтальном участке $A_{\text{тр}}=-F_{\text{тр}}s$, откуда $F_{\text{тр}}=A_{\text{тяги}}/s=1000/50=20\,\text{Н}$. Ответ: $20\,\text{Н}$.
}

% Source 354CDE: https://ege.fipi.ru/bank/questions.php?search=1&proj=BA1F39653304A5B041B656915DC36B38&qid=354CDE
\BankTask{354CDE}{
Тело массой $0{,}1\,\text{кг}$ вращается в горизонтальной плоскости на нити длиной
$1\,\text{м}$. Чему равна работа силы тяжести за один оборот вращения тела? Результат выразите в $\text{Дж}$.
}{
\smallskip\noindent\textbf{Решение от GPT}\par
При движении в горизонтальной плоскости $\Delta h=0$. Работа силы тяжести $A=-mg\Delta h=0\,\text{Дж}$.
}

% Source FB0F2e: https://ege.fipi.ru/bank/questions.php?search=1&proj=BA1F39653304A5B041B656915DC36B38&qid=FB0F2e
\BankTask{FB0F2e}{
В первой серии опытов брусок перемещают при помощи нити равномерно вверх по наклонной плоскости. Во второй серии опытов на брусок ставят груз и продолжают перемещать брусок при помощи нити равномерно вверх по той же наклонной плоскости. В обоих случаях нить параллельна наклонной плоскости.
Как изменяются при переходе от первой серии опытов ко второй модуль работы силы трения при перемещении бруска на одинаковые расстояния и коэффициент трения скольжения между бруском и плоскостью?
\begin{enumerate}
\item увеличивается
\item уменьшается
\item не изменяется
\end{enumerate}
\TaskChangeTable{Модуль работы силы трения}{Коэффициент трения скольжения}
}{
\smallskip\noindent\textbf{Решение от GPT}\par
При той же плоскости $\mu$ неизменен. После добавления груза $N=(m+M)g\cos\alpha$ больше, поэтому $|A_{\text{тр}}|=\mu Ns$ на том же пути увеличивается. Ответ: $13$.
}

\end{document}
